% Notes-only candidate. No canonical registration before complete proof audits.
\subsection{A descent-kernel element invisible on a full derived quotient}

\begin{statement}{C.46}{Open}{Candidate: parity survives beyond the third-derived quotient}
Write $F^{(0)}=F$ and $F^{(i+1)}=[F^{(i)},F^{(i)}]$, with
closed commutator subgroups. There is an actual multiplier-one full
gentle element $e$ whose normalized original conjugator belongs to
$F^{(3)}=F'''$ and for which
\[
 D(e)=1,\qquad \Delta(e)=1,\qquad
 e|_{F/F'''}=1.                                      \tag{C.46.1}
\]
The last equality is the identity on the full profinite derived quotient,
not only on a selected finite quotient. Consequently $\Delta$ does not
factor through the actual gentle action on $F/F'''$, and is linearly
independent of every family of characters factoring through that action.
In particular it is independent of the metabelian beta-defect characters
of C.30. This assertion does not give further preimages under $D$,
surjectivity of $D$ on the gentle group, or independence of all iterates
$\Delta D^n$.
\end{statement}

\begin{proof}[Candidate proof, pending complete audits]
Use the actual target modular group and its index-six pure subgroup
\[
 \Gamma=\widehat{C_2*C_3}=\langle s,r\mid s^2=r^3=1\rangle,
 \qquad \tau=sr,\qquad F=\ker(\Gamma\longrightarrow S_3).
\]
Thus the derived subgroups in this proof are those of this specified
$F$. As in C.23 and C.44, take the temporary marking
\[
 x=\tau^{-2},\qquad y=sxs,\qquad z=(xy)^{-1}.
\]
The original filling, expressed in these temporary letters, is
\[
 \varphi(x)=\tau^{-1},\qquad \varphi(y)=s,\qquad
 \varphi(z)=r.                                       \tag{C.46.2}
\]
Let $\rho=\operatorname{Int}(r)$ and
$\theta=\operatorname{Int}(s)$. Their six actual peripheral
representatives are $1,\rho,\rho^2,\theta,\rho\theta,
\rho^2\theta$, with substituted pairs
\[
 (x,y),\ (y,z),\ (z,x),\ (y,x),\ (z,y),\ (x,z).
                                                               \tag{C.46.3}
\]
Set
\[
 U=F''\cap\bigcap_\sigma\ker(\varphi\sigma:F\to\Gamma).
                                                               \tag{C.46.4}
\]
This is a closed normal subgroup of $\Gamma$. Indeed $F''$ is
characteristic in $F$, and multiplication of the six representatives
changes a representative only by an inner automorphism of $F$;
normality of each filling kernel in $F$ therefore makes their
intersection invariant under $\Gamma$. No openness is asserted.

\paragraph{An exact word in the second derived subgroup.}
With $[a,b]=aba^{-1}b^{-1}$, put
\[
 a=[x^2,y^2],\qquad b=[y^2,z^2],\qquad w=[a,b].       \tag{C.46.5}
\]
Both $a$ and $b$ lie in $F'$, so $w\in F''$. Under every one
of the six full fillings, one of the peripheral classes $x,y,z$
becomes an involution. If $x^2$ becomes one, then $a=1$;
if $y^2$ becomes one, both inner commutators are one; if $z^2$
becomes one, then $b=1$. Hence
\[
                 \varphi\sigma(w)=1\quad\text{in }\Gamma
                 \quad\text{for all six }\sigma.    \tag{C.46.6}
\]
These are full free-product identities. The executed word certificate
also reduces each of their actual substituted words to the empty
$C_2*C_3$ normal form. In particular $w\in U$.

Let $q:\Gamma\to Q_{12}=\operatorname{SL}_2(\mathbb Z/12)/\{\pm I\}$
be the original congruence quotient, with the signed pure marking
$X=U_0^2$, $Y=SU_0^2S^{-1}$ used in C.44. Exact multiplication gives
\[
\begin{array}{c|c|c|c}
 &a&b&w\\ \hline
 \bmod 3&\begin{pmatrix}1&1\\1&2\end{pmatrix}
 &\begin{pmatrix}0&2\\1&0\end{pmatrix}&-I\\[3pt]
 \bmod 4&I&I&I\\[3pt]
 \bmod 12&\begin{pmatrix}1&4\\4&5\end{pmatrix}
 &\begin{pmatrix}9&8\\4&9\end{pmatrix}&5I
\end{array}                                                   \tag{C.46.7}
\]
Thus $q(w)=[5I]$.

For completeness the entire second-derived image is a central
$2$-group. Modulo three, $X,Y$ generate
$\operatorname{SL}_2(\mathbb F_3)$, whose derived images in these
literal matrices have orders $8$ and $2$, the second being
$\{I,-I\}$. Modulo four the pure image is the abelian group of
order four generated by $X,Y$. These assertions can also be checked
by listing the finite generated group and all commutators; the
certificate does this at moduli $3,4,12$, obtaining pure-image orders
$24,4,96$ and first-derived orders $8,1,8$. Therefore
\[
       q(F'')=\{1,[5I]\},\qquad q(U)=\{1,[5I]\}.     \tag{C.46.8}
\]
The equality for $U$ follows from $w\in U$, not merely from an
upper bound on its image.

\paragraph{The full relative lift with retained pure frames.}
Apply the closed-normal-subgroup argument proved in C.44 to the
specific $U$ in (C.46.4). Write $R_{\rm nil}(U)$ for its
pronilpotent residual. The exact relative product is
\[
 \Gamma/R_{\rm nil}(U)
  \simeq\prod_{p,\,\Gamma/U}\Gamma/O^p(U).            \tag{C.46.9}
\]
This is the actual quotient of $\Gamma$: for each finite quotient
$q_E$, the image of $R_{\rm nil}(U)$ is the nilpotent residual
of $q_E(U)$. The finite scalar-cover refinements in P.3 kill this
residual because their $U$-images are central-by-nilpotent. Compactness
therefore lifts the relative automorphism and its specified three
frames without requiring $U$ open.

Here are the frames. Since $q(U)$ is a $2$-group, $q$ factors
through $\Gamma_2=\Gamma/O^2(U)$. Put
\[
 J=\operatorname{im}\bigl(\Gamma\to Q_{12}\times\Gamma/U\bigr).
\]
The map $\Gamma_2\to J$ has pro-two kernel. The element
$([5I],1)$ belongs to $J$ and centralizes the image of $r$.
Choose a preimage in $U^{(2)}$. Its conjugate of $r$ and $r$
are order-three complements in the preimage of the finite subgroup
$\langle r_J\rangle$. Coprime complement conjugacy in finite
pro-two quotients and compactness give a correction in that kernel.
Consequently there is
\[
 c\in C_{\Gamma_2}(r)\cap U^{(2)},\qquad q(c)=[5I].  \tag{C.46.10}
\]
All the other prime coordinates are one. The identity automorphism
of (C.46.9) now has the compatible prescribed frames
\[
 b_s=b_\tau=1,\qquad b_r=(c\text{ at }2,\ 1\text{ otherwise}).
                                                               \tag{C.46.11}
\]
The full relative lifting argument gives a peripheral automorphism
$\Phi$ of $\Gamma$. Its cusp frame can be made literally one by
an inner adjustment in $R_{\rm nil}(U)$. Thus $\Phi(\tau)=\tau$
exactly, while its unique pure torsion frames $A,\beta$ retain
\[
                    q(A)=[5I],\qquad q(\beta)=1.      \tag{C.46.12}
\]
The finite signed values here are values of actual retained frames.

\paragraph{Both normalizations preserve the full derived identity.}
Since $\Phi$ induces identity on $\Gamma/R_{\rm nil}(U)$, put
\[
 f=s\Phi(s)\in R_{\rm nil}(U)\subset U'\subset F''' . \tag{C.46.13}
\]
The first containment after the residual holds because every finite
abelian quotient of $U$ is nilpotent. The second holds because
$U\subset F''$. The cusp adjustment just made belongs to the same
residual, so it preserves all these containments.

The exact torsion relations give the two cycle identities, and the
restriction to $F$ is the actual normalized multiplier-one gentle
pair $x\mapsto x$, $y\mapsto f^{-1}yf$ in the temporary marking.
There is no further abelian normalization exponent: it is zero
because $f\in F'''\subset F'$. Return to the original marking
\[
             x_{\rm old}=y^{-1},\qquad y_{\rm old}=x^{-1}.
\]
Conjugation of the restriction by $f$ fixes $x_{\rm old}$ and
sends $y_{\rm old}$ to $f y_{\rm old}f^{-1}$. Thus its normalized
original conjugator is $f_{\rm old}=f^{-1}\in F'''$.
The marking change is an automorphism of $F$, so the derived
subgroup is unchanged. This proves identity on the entire quotient
$F/F'''$, as stated in (C.46.1).

Moreover $f\in U$ and every full filling kills $U$.
The same abstract original filling therefore satisfies
$\varphi(f_{\rm old})=1$ in the full $\Gamma$. The descended
modular automorphism fixes $s$ and sends $\tau$ to
$\varphi(f_{\rm old})^{-1}\tau\varphi(f_{\rm old})=\tau$.
It is the full identity, hence $D(e)=1$.

Finally $q$ kills $R_{\rm nil}(U)$, so $q\Phi=q$ and the
normalized congruence intertwiner is $V=I$ by C.12. The actual
pure frames (C.46.12) give the CM coordinates $([5I],1)$,
and hence $\Delta(e)=1$, exactly as in C.44. The original-marking
inner adjustment has $q(f)=1$, so changes none of those signed data.
This proves all three full assertions in (C.46.1).

If $\Delta$ factored through the action on $F/F'''$, it would
vanish on this $e$, a contradiction. Every metabelian beta-defect
character factors through $F/F''$ and therefore through $F/F'''$,
so all of them vanish on $e$. This proves the character distinction.

The standalone executed program
\path{code/pbc_derived_descent_kernel_word.py}
retains the freely reduced $34$-letter word, all six full substituted
normal forms, the old-marking transfer and its inverse, all signed
matrices in (C.46.7), and the complete finite-derived-image checks.
Its JSON and stdout are retained. Those finite checks support the
explicit word and frame data; the full quotient identity follows
from (C.46.13) and the actual profinite lift, not from enumeration.
\end{proof}
